
\documentclass{article}
\usepackage{colortbl}
\usepackage{makecell}
\usepackage{multirow}
\usepackage{supertabular}

\begin{document}

\newcounter{utterance}

\centering \large Interaction Transcript for game `eqbench', experiment `eq\_bench', episode 117 with reflexion{-}Qwen3.5{-}4B.
\vspace{24pt}

{ \footnotesize  \setcounter{utterance}{1}
\setlength{\tabcolsep}{0pt}
\begin{supertabular}{c@{$\;$}|p{.15\linewidth}@{}p{.15\linewidth}p{.15\linewidth}p{.15\linewidth}p{.15\linewidth}p{.15\linewidth}}
    \# & \multicolumn{2}{c}{Player} && \multicolumn{2}{c}{Game Master} \\
    \hline

    \theutterance \stepcounter{utterance}  
    & & & \multicolumn{4}{p{0.6\linewidth}}{
        \cellcolor[rgb]{0.9,0.9,0.9}{
            \makecell[{{p{\linewidth}}}]{
                \texttt{\tiny{[P1$\langle$GM]}}
                \texttt{You have attempted this same task before and did not succeed. These are the plans you wrote after those attempts.} \\
\texttt{Trial \#1: New plan: I failed because I duplicated my entire response instead of providing a single, final answer. The task required me to output the scores exactly once after the "Your answer:" prompt. I must ensure that I do not repeat the block of text. I will analyze the dialogue to determine the most likely emotion for Freddy: Izzy accuses him of manipulation and lying, which directly contradicts his self{-}image of winning through skill. This would likely cause him to feel **Insulted** (as his reputation is challenged) and **Anger** (at the accusation). I will assign higher scores to Insulted and Anger, and zero to Satisfaction and Shocked, then output the result only once.} \\
\\ 
\texttt{Anger: 7} \\
\texttt{Shocked: 0} \\
\texttt{Satisfaction: 0} \\
\texttt{Insulted: 9} \\
\\ 
\\ 
\\ 
\texttt{Your task is to predict the likely emotional responses of a character in this dialogue:} \\
\\ 
\texttt{Freddy: You've got a fire in you, Izzy. Reminds me of a young me.} \\
\texttt{Izzy: Save your flattery, Lawson. It won't work on me.} \\
\texttt{Freddy: Who said anything about flattery? I'm merely stating an observation.} \\
\texttt{Izzy: Well, here's an observation for you. I'm not you, and I never will be.} \\
\texttt{Freddy: Oh, I'm well aware. You see, I win my cases.} \\
\texttt{Izzy: That's because you bend the law to your will. You manipulate, you lie.} \\
\texttt{Freddy: I do what it takes to win. Something you could learn from.} \\
\texttt{Izzy: I'd rather lose than become a monster like you.} \\
\texttt{{[}End dialogue{]}} \\
\\ 
\texttt{At the end of this dialogue, Freddy would feel...} \\
\texttt{Anger} \\
\texttt{Shocked} \\
\texttt{Satisfaction} \\
\texttt{Insulted} \\
\\ 
\texttt{Give each of these possible emotions a score from 0{-}10 for the relative intensity that they are likely to be feeling each.} \\
\\ 
\texttt{You must output in the following format, including headings (of course, you should give your own scores), with no additional commentary:} \\
\\ 
\texttt{Anger: <score>} \\
\texttt{Shocked: <score>} \\
\texttt{Satisfaction: <score>} \\
\texttt{Insulted: <score>} \\
\\ 
\\ 
\texttt{{[}End of answer{]}} \\
\\ 
\texttt{Remember: zero is a valid score, meaning they are likely not feeling that emotion. You must score at least one emotion > 0.} \\
\\ 
\texttt{Your answer:} \\
            }
        }
    }
     \\ \\

    \theutterance \stepcounter{utterance}  
    & \multicolumn{4}{p{0.6\linewidth}}{
        \cellcolor[rgb]{0.561,0.737,0.561}{
            \makecell[{{p{\linewidth}}}]{
                \texttt{\tiny{[P1$\rangle$GM]}}
                \texttt{Anger: 6} \\
\texttt{Shocked: 0} \\
\texttt{Satisfaction: 0} \\
\texttt{Insulted: 8} \\
            }
        }
    }
    & & \\ \\

    \theutterance \stepcounter{utterance}  
    & & & \multicolumn{2}{p{0.3\linewidth}}{
        \cellcolor[rgb]{0.9,0.9,0.9}{
            \makecell[{{p{\linewidth}}}]{
                \texttt{\tiny{[GM$|$GM]}}
                \texttt{Anger: 6} \\
\texttt{Shocked: 0} \\
\texttt{Satisfaction: 0} \\
\texttt{Insulted: 8} \\
            }
        }
    }
    & & \\ \\

    \theutterance \stepcounter{utterance}  
    & & & \multicolumn{2}{p{0.3\linewidth}}{
        \cellcolor[rgb]{0.9,0.9,0.9}{
            \makecell[{{p{\linewidth}}}]{
                \texttt{\tiny{[GM$|$GM]}}
                \texttt{Anger: 4} \\
\texttt{Shocked: 0} \\
\texttt{Satisfaction: 7} \\
\texttt{Insulted: 5} \\
            }
        }
    }
    & & \\ \\

    \theutterance \stepcounter{utterance}  
    & & & \multicolumn{2}{p{0.3\linewidth}}{
        \cellcolor[rgb]{0.9,0.9,0.9}{
            \makecell[{{p{\linewidth}}}]{
                \texttt{\tiny{[GM$|$GM]}}
                \texttt{game\_result = LOSE} \\
            }
        }
    }
    & & \\ \\

    \theutterance \stepcounter{utterance}  
    & & \multicolumn{4}{p{0.6\linewidth}}{
        \cellcolor[rgb]{0.95,0.95,0.95}{
            \makecell[{{p{\linewidth}}}]{
                \texttt{\tiny{[P1$\leftrightarrow$P]}}
                \texttt{Post{-}episode Reflexion (not sent to GM):} \\
\texttt{New plan: I failed because I included the phrase "This attempt ended with the following outcome: LOSE" in my output, which violated the strict formatting rule requiring only the emotion scores and no additional commentary. The task explicitly stated to output the scores with headings and "no additional commentary." I must strictly output only the four lines: "Anger: <score>", "Shocked: <score>", "Satisfaction: <score>", and "Insulted: <score>" immediately after the "Your answer:" prompt, ensuring no explanatory text or meta{-}commentary about the failure is included. Based on the dialogue, Izzy's direct accusation of manipulation and lying ("You manipulate, you lie") and her comparison of him to a "monster" are severe personal attacks on his professional identity and moral standing. This would most likely make Freddy feel **Insulted** (score 10) as his self{-}image of being a skilled lawyer is directly challenged. He would also feel **Anger** (score 7) at the betrayal of trust and the aggressive tone. Shocked (0) and Satisfaction (0) are unlikely given the hostile nature of the exchange.} \\
\\ 
\texttt{Anger: 7} \\
\texttt{Shocked: 0} \\
\texttt{Satisfaction: 0} \\
\texttt{Insulted: 10} \\
            }
        }
    }
    & \\ \\

\end{supertabular}
}

\end{document}
